\subsection{TOPOLOGICAL DIVISION RINGS AND FIELDS}

In what follows, and in Chapters IV and V, if K is a division ring we shall denote by \(K^{*}\) the multiplicative group of non-zero elements of K.

DEFINITION 4. A topological division ring is a set K carrying a division ring structure and a topology compatible with the ring structure of K, and satisfying in addition the following axiom:

(KT) The mapping \(x \to x^{-1}\) of \(\mathbf{K}^*\) into \(\mathbf{K}^*\) is continuous.

A commutative topological division ring is called a topological field.

A division ring structure and a topology on a set K are said to be compatible if the corresponding ring structure and the topology are compatible and if in addition axiom (KT) is satisfied.

Examples. 1) On any division ring K, the discrete topology is compatible with the division ring structure. A topological division ring whose topology is discrete is called a discrete division ring.

* 2) The topology of the rational line Q (resp. the real line R) is compatible with the field structure of R (resp. R) (see Chapter IV, § 3).*

Definition 4 shows that, if K is a topological division ring, then the topology induced by that of K on the multiplicative group \(K^{*}\) is compatible with the group structure of \(K^{*}\) .

If \(a \neq 0\) , the homotheties \(x \rightarrow ax\) and \(x \rightarrow xa\) are homeomorphisms of K onto itself; and so is the mapping \(x \rightarrow ax + b\) for all \(b \in K\) . Note that the homotheties \(x \rightarrow ax\) and \(x \rightarrow xa\) are automorphisms of the (topological) additive group of K if \(a \neq 0\) . If V is any neighbourhood of o in K, it follows therefore that aV and Va are neighbourhoods of o for all \(a \neq 0\) .

Let X be a topological space, and let f be a mapping of X into a topological division ring K. If f is continuous at a point \(x_{0} \in X\) and if \(f(x_{0}) \neq 0\) , then \(f^{-1}\) is continuous at \(x_{0}\) . In particular, if K is a topological field, then every rational function in n variables with coefficients in K is continuous at every point of \(K^{n}\) where its denominator does not vanish.

Likewise, if \(f\) is a mapping of a set \(\mathbf{X}\) , filtered by a filter \(\mathfrak{F}\) , into a Hausdorff topological division ring \(\mathbf{K}\) , and if \(\lim_{\mathfrak{F}} f\) exists and is not \(o\) , then \(\lim_{\mathfrak{F}} f^{-1}\) exists and we have

\[
\lim _ {\mathfrak {F}} f ^ {- 1} = (\lim _ {\mathfrak {F}} f) ^ {- 1}.\tag{4}
\]

If H is a division subring of a topological division ring K, then the topology induced on H by the topology of K is compatible with the division ring structure of H. The structure of a topological division ring thus defined on H is said to be induced by that of K. Furthermore, \(\overline{H}\) is also a division subring of K (the proof is analogous to that of Proposition 5).

In a topological division ring K, the closure of the set \(\{0\}\) is a two-sided ideal, by Proposition 5, and hence must be either \(\{0\}\) or K. In other words, if the topology of K is not the coarsest topology (Chapter I, § 2, no. 2) then it is Hausdorff (§ 1, no. 2, Proposition 2).

